We conclude with two remarks on the group \(\Gamma_0(R)\).

\begin{proposition}\label{sga3:XXI:3.5.5}
Let \(G\) be an abelian group, and let \(f:R\to G\) be a map
satisfying:
\begin{enumerate}
\item[(i)] if \(\alpha\in R\), then
  \(f(-\alpha)=-f(\alpha)\);
\item[(ii)] if \(\alpha,\beta,\alpha+\beta\in R\), then
  \[
   f(\alpha+\beta)=f(\alpha)+f(\beta).
  \]
\end{enumerate}
There exists a unique group homomorphism
\[
 \overline f:\Gamma_0(R)\longrightarrow G
\]
such that \(\overline f(\alpha)=f(\alpha)\) for every
\(\alpha\in R\).
\end{proposition}

\begin{proof}
Let \(\Delta\) be a system of simple roots of \(R\), and write
\[
 \beta=\sum_{\alpha\in\Delta}a(\alpha)\alpha.
\]
Lemma~\ref{sga3:XXI:3.2.16} gives at once
\[
 f(\beta)=\sum_{\alpha\in\Delta}a(\alpha)f(\alpha).
\]
Now \(\Delta\) is a basis of \(\Gamma_0(R)\).
\end{proof}

\begin{proposition}\label{sga3:XXI:3.5.6}
Let \(\Delta\) be a system of simple roots. There exists on
\(\Gamma_0(R)\) a totally ordered group structure such that the roots
\(>0\) are the elements of \(\mathscr{P}(\Delta)\), and such that
\[
 \alpha\longmapsto\operatorname{ord}_\Delta(\alpha)
\]
is increasing.
\end{proposition}

\begin{proof}
Let
\[
 \alpha_1,\alpha_2,\ldots,\alpha_n,
 \qquad n=\rgss(\cR),
\]
be the elements of \(\Delta\). Every \(x\in\Gamma_0(R)\) has a
decomposition
\[
 x=\sum_{i=1}^{n}m_i(x)\alpha_i.
\]
Take the lexicographic order with respect to the functions
\[
 \sum_i m_i,\quad m_n,\quad m_{n-1},\quad\ldots,\quad m_2.
\]
\end{proof}

\begin{remark}\label{sga3:XXI:3.5.7}
The first roots, in order, are
\[
 \alpha_1,\alpha_2,\ldots,\alpha_n;
\]
then come, when they are roots,
\[
 2\alpha_1,\quad
 \alpha_1+\alpha_2,\quad
 \alpha_1+\alpha_3,\quad\ldots.
\]
\end{remark}

\subsection{Weyl chambers}
\label{sga3:XXI:3.6}

\begin{lemma}\label{sga3:XXI:3.6.1}
Let \(V\) be a finite-dimensional real vector
space.\footnote{Editors' note: The original \(\mathbb{Q}\) has been
replaced by \(\mathbb{R}\).} Let \(f_i\) be independent linear forms.
Put
\[
 C=\{x\in V\mid f_i(x)>0\}.
\]
Then \(C\) is a maximal convex subset of
\[
 X=V\setminus\bigcup_i f_i^{-1}(0).
\]
\end{lemma}

\begin{proof}
Trivial.
\end{proof}
